\documentclass[5pt, a4paper]{article}

\usepackage[utf8]{inputenc}
\usepackage[T1]{fontenc}
\usepackage{textcomp}
\usepackage{amsmath, amssymb}
\usepackage{multicol}
\usepackage{proof}
\usepackage{enumitem}
\usepackage{fancyhdr}
\usepackage[margin=1in]{geometry}
% figure support
\usepackage{import}
\usepackage{xifthen}
\pdfminorversion=7

\usepackage[most]{tcolorbox}
\usepackage{pdfpages}
\usepackage{transparent}
\newcommand{\incfig}[1]{%
    \def\svgwidth{\columnwidth}
    \import{./figures/}{#1.pdf_tex}
}

\tcbset{
colback=gray!5,
colframe=blue!50!white,
coltitle=black,
arc=0pt,
boxrule=-0.1mm,
fonttitle=\bfseries,
coltitle=black,
left=2.5mm,
leftrule=1mm,
right=3.5mm,
colbacktitle=black!10!white,
toptitle=0.75mm,
bottomtitle=0.5mm,
}
\begin{document}
    \section{Least Square Problems} 
    \subsection*{Introduction}
    Suppose that you have a system $Ax = b$, where $A$ is a $n\times m$ matrix
    where $n > m$. In such a system, we usually say that the system is
    overdetermined, and a solution wouldn't exist. Often times when presented
    with a problem like this, such as fitting a line over a couple of data
    points (elaborated further in later chapters concering interpolation,) we
    would like to find some solution that best fit our data.

    Let's take a look at our previous example $Ax = b$. Suppose for
    $\mathbb{R}^{2}$we write them out as such:  \[
        \begin{pmatrix} 
            a_{11} & a_{12}  \\
            a_{21} & a_{22}  \\
            a_{31} & a_{32}  \\
            a_{41} & a_{42} \\
        \end{pmatrix}    
        \begin{pmatrix} 
            x_1 \\ 
            x_2
        \end{pmatrix} 
        = \begin{pmatrix} y_1 \\ y_2 \\ y_3 \\y_4 \end{pmatrix} 
    .\] 
     Here, we couldn't find an exact solution to the problem. But, we can come
     up with some solution that minimize the 'error'. Such a problem is called a
     least square problem. 

     To understand how we could solve such a problem, we'll analyze $Ax = b$
     once more. If we think of $A$ as as some plane in $\mathbb{R}^{2}$, we can
     think of the existance of a solution for $Ax = b$ as that  $b$ lies
     somewhere within the span of $A$. Likewise, if there is no solution, $b$
     lise somewhere outside the span of $A$. Notice, however, that we can find
     some vector $y$ within the span of $A$ ($A\hat{x} = y$) such that $b = y + e$ for some
     vector $e$ taht satisfies  $A^{T}e = 0$ ($e$ is perpendicular to  the plane
     of $A$). As such, we can write the following: \[
     0 = A^{T}e = A^{T}(b - y) = A^{T}(b - A \hat{x})
     .\] 
     \[
     A^{T}b = A^{T}A \hat{x}
     .\] 
     We can interpret the solution as finding the vector $\hat{x}$ that
     minimizes $e$ (as it so happens for $\mathbb{R}^{2}$, the minimum is just
     the projection of $b$ unto the plane.)

     Notice the aforementioned expression \[
     A^{T}b = A^{T}A\hat{x}
     .\] 
     We can rearrange it into the following: \[
         (A^{T}A)^{-1}A^{T}b = \hat{x}
     .\] 
     We define \[
        A^{+} = (A^{T}A)^{-1}A^{T}
     .\] 
     Where $A^{+}$ is the pseudoinverse of $A$. It can be shown that for any
     $n\times m$ matrix $A$, $A$ has a pseudoinverse that is positive.

     \subsection*{QR Factorization}
     A QR factorization is simply decomposing a vector into $A = QR$, where $Q$ 
     is a orthogonal matrix $Q = Q^{T}$ and $R$ is an upper triangular matrix.
     We could interpret this as taking an orthonormal basis $Q$ and constructing
     $A$ using $R$. Generally, it is preffered over LU as it is more stable and
     does not require pivoting as LU does.

    When solving a $QR$ factorization, we have a couple of options: 
     \begin{enumerate}
        \item Gramm-Schmidtt Algorithm
        \item Householder Reflections 
        \item Givens Rotation
        \item Singular Value Decomposition (SVD)
    \end{enumerate}
    Overall, Gramm-Schmidtt is not preffered as it quite slow. We'll talk about
    the first three in this chapter, and $SVD$ in the later chapters concerning
    eigenvalue problems.

    \begin{tcolorbox}[title=Gramm-Schmidtt]
         The idea behind Gramm-Schmidtt is to construct a set of orthonormal
         vectors by considering each columns of the matrix $A$, and iteratively
         subtracting it by it's projections on each previous columns and
         normalizing it, while storing the dot product of each vectors in an
         upper triangular matrix.
    \end{tcolorbox}

    As one might expect, when doing the calculation by relying on previous
    results of computing each row of $Q$, error would be magnified. As such,
    when calculating a QR decomposition using this method, we would design the
    algorithm to iterate by rows instead.
    \subsection*{Rotations and Reflections}

    Let's start with Householder Relfections. In short, a Householder reflection
    is an algorithm that, for each vector column in a matrix, try to reflect it
    such that it makes it orthogonal, then storing the reflector in a upper
    triangular matrix.   

    Before going further, let's  discuss as to what it means for a vector to be
    reflected. Suppose some vector $\hat{b} \in \mathbb{R}^2$ that we wish to
    reflect in such a way that it is in the span of $e_x$. To do so, the easiet
    way is to find some line $\hat{w}$ that bisects the angle that $\hat{x}$
    makes with the $x$ axis, and subtract two times the vector $\hat{v}$ that is
    perpindicular to that line from $\hat{b}$. So esentially, reflecting is
    subtracting two times the perpendicular element of our choosen reflector,
    from $\hat{b}$.

    In a Householder Reflection, we wish to do a relfection on each vector
    column of the matrix such that the entries below the pivot is zero'd out. To
    compute such a reflection, we follow \[
        H_i = I - \frac{\hat{v}\hat{v}^{T}}{\hat{v}^{T}\hat{v}}
    .\] 
    We'll then solve for $\hat{v}$ such that the expression \[
        H\hat{a}_i = \alpha e_i
    s\] 
    is statisfied. We'll repeat for each collumn until we have $R$, and we set
     $Q = H_1H_2\ldots H_n$.

     Even though this method is generally faster, for sparse matrices, or those
     with special structure, a Givens Rotation might be easier to perform.
     Contrary to a Householder Reflection, which seeks to zero out every single
     entry under the pivot in one go, a Givens rotation instead calculates a
     raotaion matrix that would zero out each individual entries. 

     Rotating a vector is another way to achieve the same effect, as they both
     preserve the norm. Going back to the first example, a vector $\hat{b}$ can
     be moved unto the span of $e_x$ by rotating it, instead of a reflection. To
     do so, we multiply the vector $\hat{b}$ by the rotation matrix \[
     R(\theta) = \begin{pmatrix} 
         \cos \theta &  \sin \theta \\ -\sin \theta & \cos \theta \\
     \end{pmatrix} 
     .\] 
     To do construct the matrix, we'll need to compute the sine and cosine of
     the rotation. Suppose that $\hat{b} = \begin{pmatrix} x, y \end{pmatrix} $.
     Then, we have \[
      \sin \theta = \frac{y}{\sqrt{x^2+y^2} }, \quad \cos \theta =
      \frac{x}{\sqrt{x^2+y^2} }
     .\] 
     For larger vectors, a intuitive way we can think of it is as 'adjusting' a
     line such that it lies only in one axis, and doing so by moving it one axis
     at a time.

     In dering the QR decomposition, we do this  rotation for each entry below
     the pivot, until we arrive at $R$, whilst accumulating the  $Q$

    

\end{document}
